\section*{Kapitel 8 --- SPI och transaktionsprotokollet}
\addcontentsline{toc}{section}{Kapitel 8 --- SPI och transaktionsprotokollet}

\solution{ex:8:1}{Avkoda}
Kommandobyten är \code{W} i bit 7, nollor i bitarna 6--4, och registerindexet i bitarna 3--0.
\ans{a} \code{0x81} = skrivbit satt, index 1: en \textbf{skrivning till \code{TX_ID}} med värdet
\code{0x000007FF} --- elva bitar ettor, alltså det största värde registret kan hålla.
\ans{b} \code{0x85} = skrivning, index 5: en \textbf{skrivning till \code{TX_SEND}} med
\code{0x1}, som \emph{utlöser en sändning}. Registret lagrar ingenting.
\ans{c} \code{0x8B} = skrivning, index 11: en \textbf{skrivning till \code{ERROR_FLAGS}} med
\code{0x0}, som nollställer felbiten.
\ans{d} \code{0xC2} är \code{1100 0010}: skrivbiten satt, men \emph{också} bit 6, som är
reserverad och måste vara noll. Kommandot är ogiltigt --- ingen skrivning verkställs, ingen
läsning låses --- och slaven konsumerar ändå alla fem bytes innan den återgår till vila.

De två ovanliga är \textbf{c} och \textbf{d}. I c nollställer man genom att skriva \emph{noll},
vilket är motsatt regel mot \code{TX_SEND} och \code{RX_ACK} som kräver bit 0 satt. I d är
kommandot ogiltigt och slaven rapporterar det \emph{inte}: det finns ingen felkanal på det här
lagret, så en drivrutin som skickar skräp får tyst nonsens tillbaka.

\solution{ex:8:2}{Koda}
\ans{a} Läsning har bit 7 låg, så kommandobyten är bara indexet; \code{RX_DLC} har index 7 och de
fyra databyten är dummybytes: \code{07 00 00 00 00}
\ans{b} \code{82 00 00 00 04}
\ans{c} \code{85 00 00 00 01}
\ans{d} \code{8A 00 00 00 01}
\ans{e} \code{8B 00 00 00 00}

Skillnaden mellan d och e är regeln för vad som räknas som en giltig skrivning. \code{RX_ACK}
kräver att \textbf{bit 0 är satt} --- en nolla gör ingenting. \code{ERROR_FLAGS} kräver att
\textbf{hela ordet är noll} --- \code{0x1} gör ingenting. Två register bredvid varandra i samma
karta, med motsatta regler, och att blanda ihop dem är en av projektets vanligaste buggar.

\solution{ex:8:3}{Vad kommer tillbaka}
\ans{a} Ingen. Under en skrivning är allt slaven driver på \code{MISO} meningslöst.
\ans{b} Fyra: de sista fyra. Byten som klockas in \emph{under kommandobyten} är rester från förra
gången, eftersom slaven laddar sin svarsskiftare först i slutet av kommandobyten.
\ans{c} Returvärdet blir skiftat en byte: den första riktiga databyten hamnar i bitarna 23--16
i stället för 31--24, och den sista databyten faller bort helt. I felsökningen ser det ut som att
varje registervärde är $256$ gånger för litet, eller --- för \code{STATUS} --- som att biten man
letar efter aldrig sätts. Det är samma symptom som när \code{SPI0.DATA} inte läses i
\code{transfer()}, vilket gör de två lätta att förväxla.

\solution{ex:8:4}{\texorpdfstring{\code{SS}}{SS} och avbrott}
\ans{a} Transaktionen överges helt, utan sidoeffekter: \code{TX_ID} ändras inte, ingen trigger
går, och slaven återgår till vila redo för nästa fallande \code{SS}-flank. Det är vad som gör
slaven självläkande efter en glitch.
\ans{b} Slaven tar de fem första byten som en transaktion och \textbf{kastar de fem följande}. En
transaktion börjar bara på den första byten efter att \code{SS} fallit, så den sjätte byten
tolkas aldrig som en ny kommandobyte.

Det farliga är att det misslyckas \emph{tyst}. Vid bring-upen ser det ut som att varannan
registerskrivning försvinner: den första kommer fram, den andra inte, och ingenting rapporterar
fel. Man misstänker kopplingen, klockfrekvensen och drivrutinens bitpackning --- långt innan man
misstänker \code{SS}.
\ans{c} \code{0x91} är \code{1001 0001}: skrivbiten satt, index 1, men bit 4 är satt och
reserverad. Kommandot är ogiltigt och ignoreras helt. Det ska förhindras redan när kommandobyten
byggs: \code{CmdWrite | index} med ett index i intervallet 0--12 kan aldrig sätta bitarna 6--4,
vilket är skälet till att indexet är en namngiven konstant och inte ett tal från en beräkning.

\solution{ex:8:5}{Tidsbudget}
\ans{a} Sex transaktioner --- \code{STATUS}, \code{RX_ID}, \code{RX_DLC}, två dataregister och
\code{RX_ACK} --- alltså omkring $6 \times 45 = 270$~µs.
\ans{b} En maximal frame är omkring 130 bitar, alltså cirka \textbf{130~µs} vid 1~Mbit/s.
\ans{c} De skriver över den föregående. Kontrollern håller exakt en mottagen frame, och den nyaste
vinner.

Det är \textbf{inte} din bugg: det är hårdvarans dokumenterade begränsning, och den är avsiktlig.
Du kan minska risken genom att polla ofta och inte göra något långsamt mellan anropen, men du kan
inte komma förbi att en \code{receive()} tar dubbelt så lång tid som framen den hämtar. Vid
tillräcklig trafikintensitet tappas frames oavsett vad drivrutinen gör.
